\chapter{Introdução}
\label{cap:introducao}

% Fontes no wiki: Thesis Goal, Synthesis, LLM Integration Patterns.

\section{Contexto}
\label{sec:contexto}

Sistemas embarcados de missão crítica, como controle aeroespacial, frenagem automotiva e \textit{drivers} de dispositivos, reagem ao hardware por meio de
interrupções. Quando um periférico sinaliza um evento, o processador suspende o
código em execução e desvia para uma rotina de tratamento de interrupção
(\textit{Interrupt Service Routine}, \isr). Ao término da rotina, a execução
retorna ao ponto em que foi interrompida.

Esse mecanismo introduz concorrência mesmo em programas sem \textit{threads}:
uma \isr\ pode ser disparada entre duas instruções quaisquer do código
principal, e ambos podem ler e escrever as mesmas variáveis globais. Se o
programador não proteger essas variáveis, o resultado do programa passa a
depender do instante exato em que a interrupção ocorreu.
\citeonline{sdracer} relatam um caso concreto em uma versão do uCLinux, em que
um \textit{driver} serial e um tratador de interrupção compartilham a linha de
comunicação e, num caminho de execução raro, transmitem dados
simultaneamente.

\section{O problema}
\label{sec:problema}

Este trabalho trata de duas classes de defeito de concorrência em programas
orientados a interrupções. Um \textit{data race} ocorre quando dois acessos à
mesma posição de memória, em fluxos diferentes e sendo ao menos um deles uma
escrita, podem acontecer sem ordem definida entre si. Uma violação de
atomicidade ocorre quando dois acessos de um mesmo fluxo, que o programador
supunha indivisíveis, são separados por um acesso de uma \isr\ à mesma
variável, de modo que o resultado não corresponde a nenhuma execução em
sequência. Os dois conceitos são definidos com precisão no
Capítulo~\ref{cap:fundamentacao}.

Esses defeitos são difíceis de detectar, isolar e corrigir, porque dependem da intercalação das execuções \cite{sdracer}: uma interrupção precisa chegar num intervalo de poucas instruções, e uma bateria de testes pode passar em todas as execuções observadas sem nunca exercitar esse intervalo. Por isso a literatura recorre a técnicas que raciocinam
sobre todas as execuções possíveis: análise estática, execução simbólica e
verificação de modelos.

\section{A lacuna}
\label{sec:lacuna}

A revisão da literatura (Capítulo~\ref{cap:relacionados}) identifica três
lacunas que motivam este trabalho.

As ferramentas priorizam precisão e medem mal o recall. Três das quatro
ferramentas estudadas relatam zero ou quase zero falsos positivos nos próprios
benchmarks, enquanto a ausência de falsos negativos é afirmada por inspeção
manual ou dentro de um limite de desenrolamento
\cite{sdracer,intrace,nichecker}. \citeonline{bmc4av} afirmam que uma ferramenta
com 100\% de precisão encontrou apenas 39,4\% das violações existentes na mesma
suíte. Como se discute na Seção~\ref{sec:rel-comparacao}, essa afirmação depende também de como os defeitos são contados. Essa cultura de avaliação tem
origem identificável: o benchmark que todos usam foi criado para uma
competição que, nos casos simples, penalizava um falso positivo em 3 pontos e premiava uma
detecção com apenas 2 \cite{racebench}.

Nenhuma ferramenta trata as duas classes de defeito, e só a mais antiga
repara. As ferramentas de \textit{data race} e as de violação de atomicidade
formam linhas separadas, e apenas o SDRacer, de 2020, propõe correções
\cite{sdracer}.

Modelos de linguagem já ajudam a análise estática, mas nunca em
concorrência. Trabalhos recentes usam modelos de linguagem de grande escala
(\llm) para decidir quais alertas de um analisador são reais
\cite{llift,chatgpt-static,reducing-fp}, inferir especificações
\cite{iris,adataint} e propor correções \cite{skipanalyzer,sastgenius}. Todos
tratam defeitos sequenciais: fluxo de dados contaminados, desreferência de
ponteiro nulo, vazamento de recursos, uso antes da inicialização.
\citeonline{llift} citam explicitamente a concorrência como algo que a análise
estática modela mal, sem tratá-la. A interseção entre as duas literaturas está
vazia, e é nela que este trabalho se situa.

\section{Objetivos}
\label{sec:objetivos}

\subsection{Objetivo geral}

Construir uma ferramenta que encontre \textit{data races} e violações de
atomicidade em código C embarcado orientado a interrupções sem falsos
negativos em relação a um conjunto explícito de premissas, aceitando falsos
positivos, e que feche o ciclo com um modelo de linguagem: detectar,
contextualizar, triar, reparar e re-verificar.

\subsection{Objetivos específicos}

\begin{enumerate}
  \item Detectar, numa única análise estática, os candidatos das duas classes de
    defeito, com recall de 100\% sobre o gabarito do Racebench.
  \item Emitir, para cada candidato, um registro de contexto com as evidências de
    que um revisor precisa: acessos, fluxos, prioridades, caminhos de chamada e
    estado de mascaramento.
  \item Descartar apenas candidatos que um solver prove impossíveis.
  \item Usar um \llm\ para ordenar e explicar os candidatos restantes, reduzindo
    o esforço de revisão sem nunca descartar um candidato.
  \item Medir a contribuição de cada técnica de \textit{prompt}, numa ablação no
    formato de \citeonline{llift}.
  \item Propor correções específicas de interrupção e verificar que o defeito
    desapareceu sem introduzir outros.
\end{enumerate}

\section{Justificativa}
\label{sec:justificativa}

Do ponto de vista prático, um defeito de concorrência que escapa aos testes se manifesta em campo, num sistema de controle físico. Uma análise que não perde defeitos é o que permite
afirmar que um programa foi verificado. O preço dessa garantia são os falsos
alarmes: na ferramenta IntRace, a primeira etapa gera cerca de 208 candidatos
por programa real \cite{intrace}. Revisá-los um a um é caro. A etapa com \llm\
proposta aqui ataca exatamente esse custo, ordenando os candidatos de modo que
o revisor encontre os defeitos reais lendo o mínimo possível, sem que nenhum
defeito saia do relatório.

Do ponto de vista científico, o trabalho contribui em três frentes: é a
primeira aplicação de triagem por \llm\ a concorrência com interrupções; unifica
num só \textit{front end} as duas classes de defeito que a literatura trata
separadamente; e explicita as premissas sob as quais a ausência de falsos
negativos vale, o que nenhum dos trabalhos revisados faz.

\section{Organização do documento}
\label{sec:organizacao}

O Capítulo~\ref{cap:fundamentacao} apresenta os conceitos necessários:
interrupções, as duas classes de defeito, as técnicas de análise e as métricas
de avaliação. O Capítulo~\ref{cap:relacionados} apresenta e compara os
trabalhos relacionados nas duas literaturas. O Capítulo~\ref{cap:proposta}
descreve a ferramenta proposta, e o Capítulo~\ref{cap:metodologia}, como ela
será avaliada.
\ifparcial
O Capítulo~\ref{cap:plano} apresenta o plano de trabalho.
\else
Os Capítulos~\ref{cap:implementacao}, \ref{cap:avaliacao} e
\ref{cap:conclusao} apresentam a implementação, a avaliação e as conclusões. O
Apêndice~\ref{ap:contratos} detalha os contratos de dados entre as etapas da
ferramenta.
\fi
